% GENERATED FILE. Edit canonical lessons, then run npm run generate:books.
% canonical-source-hash: fnv1a64:93119778f9e6f198
% canonical-lessons: UR-C461-lihaz, UR-C461-shaista, UR-C461-tahzib, UR-C461-akhlaq, UR-C461-ehtiram-karna

\chapter{Manners}
\label{ch:ur-manners}
\hlchaptermodality{461}

\begin{chapteropening}
\textbf{By the end of this chapter:} \emph{I can say whether someone is polite or rude.}

Complete the last lesson of chapter 461: I can say whether someone is polite or rude.
\end{chapteropening}

\section[\texorpdfstring{\ur{لحاظ} (lihāz)}{لحاظ (lihāz)}]{\ur{لحاظ} (lihāz) — consideration; regard (for others)}

\label{lesson:UR-C461-lihaz}



\begin{quote}
Before the new one: say the Urdu for to fight, then the Urdu for valiant.
\end{quote}

\subsection*{You'll want to know: \ur{لحاظ}}
\textbf{\ur{لحاظ}} — \emph{lihāz} — \textquotedblleft{}consideration; regard (for others)\textquotedblright{}.

\begin{quote}
\ur{ہمیں} \ur{بڑوں} \ur{کا} \ur{لحاظ} \ur{رکھنا} \ur{چاہیے۔}
\end{quote}

\begin{quote}
\emph{hameṅ baṛoṅ kā lihāz rakhnā chāhiye.}
\end{quote}

\begin{quote}
We should show consideration to our elders.
\end{quote}

Masculine.

\begin{cousinweb}
\textbf{\ur{لحاظ}} (\emph{lihāz}) is usually traced, through Persian, to Arabic \emph{laḥāẓ}, \textquotedblleft{}a glance\textquotedblright{}, and so \textquotedblleft{}regard\textquotedblright{} for someone.
\end{cousinweb}

\begin{grammarlens}[title={What you've built}]
One more word for describing what you lived and felt.
\end{grammarlens}

\subsection*{Your turn}
\begin{itemize}
\raggedright
  \item \emph{Say it:} \emph{lihāz}
  \item \emph{Say it:} \emph{lihāz}, once more
  \item \emph{Say it:} \emph{hameṅ baṛoṅ kā lihāz rakhnā chāhiye.} — We should show consideration to our elders.
  \item \emph{Recall:} say the Urdu for to fight, then the Urdu for valiant, then say \emph{lihāz} again
\end{itemize}

\begin{tcolorbox}[breakable,colback=teal!4,colframe=teal!35!black,title={Before you move on}]
What does \ur{لحاظ} mean? (\textquotedblleft{}Consideration; regard (for others)\textquotedblright{}.) Say it once more.
\end{tcolorbox}
\section[\texorpdfstring{\ur{شائستہ} (shāistā)}{شائستہ (shāistā)}]{\ur{شائستہ} (shāistā) — refined; courteous}

\label{lesson:UR-C461-shaista}



\begin{quote}
Before the new one: say the Urdu for valiant, then the Urdu for consideration.
\end{quote}

\subsection*{You'll want to know: \ur{شائستہ}}
\textbf{\ur{شائستہ}} — \emph{shāistā} — \textquotedblleft{}refined; courteous\textquotedblright{}.

\begin{quote}
\ur{وہ} \ur{بڑی} \ur{شائستہ} \ur{لڑکی} \ur{ہے۔}
\end{quote}

\begin{quote}
\emph{voh baṛī shāistā laṛkī hai.}
\end{quote}

\begin{quote}
She is a very courteous girl.
\end{quote}

\begin{cousinweb}
\textbf{\ur{شائستہ}} (\emph{shāistā}) is usually traced to Persian \emph{shāyista}, \textquotedblleft{}fitting, worthy\textquotedblright{}.
\end{cousinweb}

\begin{grammarlens}[title={What you've built}]
One more word for describing what you lived and felt.
\end{grammarlens}

\subsection*{Your turn}
\begin{itemize}
\raggedright
  \item \emph{Say it:} \emph{shāistā}
  \item \emph{Say it:} \emph{shāistā}, once more
  \item \emph{Say it:} \emph{voh baṛī shāistā laṛkī hai.} — She is a very courteous girl.
  \item \emph{Recall:} say the Urdu for valiant, then the Urdu for consideration, then say \emph{shāistā} again
\end{itemize}

\begin{tcolorbox}[breakable,colback=teal!4,colframe=teal!35!black,title={Before you move on}]
What does \ur{شائستہ} mean? (\textquotedblleft{}Refined; courteous\textquotedblright{}.) Say it once more.
\end{tcolorbox}
\section[\texorpdfstring{\ur{تہذیب} (tahzīb)}{تہذیب (tahzīb)}]{\ur{تہذیب} (tahzīb) — refinement; culture}

\label{lesson:UR-C461-tahzib}



\begin{quote}
Before the new one: say the Urdu for consideration, then the Urdu for courteous.
\end{quote}

\subsection*{You'll want to know: \ur{تہذیب}}
\textbf{\ur{تہذیب}} — \emph{tahzīb} — \textquotedblleft{}refinement; culture\textquotedblright{}.

\begin{quote}
\ur{اس} \ur{کی} \ur{باتوں} \ur{میں} \ur{تہذیب} \ur{ہے۔}
\end{quote}

\begin{quote}
\emph{us kī bātoṅ meṅ tahzīb hai.}
\end{quote}

\begin{quote}
There is refinement in the way she speaks.
\end{quote}

Feminine.

\begin{cousinweb}
\textbf{\ur{تہذیب}} (\emph{tahzīb}) is usually traced, through Persian, to Arabic \emph{tahdhīb}, \textquotedblleft{}refining, polishing\textquotedblright{}.
\end{cousinweb}

\begin{grammarlens}[title={What you've built}]
One more word for describing what you lived and felt.
\end{grammarlens}

\subsection*{Your turn}
\begin{itemize}
\raggedright
  \item \emph{Say it:} \emph{tahzīb}
  \item \emph{Say it:} \emph{tahzīb}, once more
  \item \emph{Say it:} \emph{us kī bātoṅ meṅ tahzīb hai.} — There is refinement in the way she speaks.
  \item \emph{Recall:} say the Urdu for consideration, then the Urdu for courteous, then say \emph{tahzīb} again
\end{itemize}

\begin{tcolorbox}[breakable,colback=teal!4,colframe=teal!35!black,title={Before you move on}]
What does \ur{تہذیب} mean? (\textquotedblleft{}Refinement; culture\textquotedblright{}.) Say it once more.
\end{tcolorbox}
\section[\texorpdfstring{\ur{اخلاق} (akhlāq)}{اخلاق (akhlāq)}]{\ur{اخلاق} (akhlāq) — manners; moral conduct}

\label{lesson:UR-C461-akhlaq}



\begin{quote}
Before the new one: say the Urdu for courteous, then the Urdu for refinement.
\end{quote}

\subsection*{You'll want to know: \ur{اخلاق}}
\textbf{\ur{اخلاق}} — \emph{akhlāq} — \textquotedblleft{}manners; moral conduct\textquotedblright{}.

\begin{quote}
\ur{اس} \ur{کے} \ur{اخلاق} \ur{بڑے} \ur{اچھے} \ur{ہیں۔}
\end{quote}

\begin{quote}
\emph{us ke akhlāq baṛe achchhe haiṅ.}
\end{quote}

\begin{quote}
He has very good manners.
\end{quote}

Masculine, used in the plural.

\begin{cousinweb}
\textbf{\ur{اخلاق}} (\emph{akhlāq}) is usually traced, through Persian, to Arabic \emph{akhlāq}, \textquotedblleft{}traits of character\textquotedblright{}.
\end{cousinweb}

\begin{grammarlens}[title={What you've built}]
One more word for describing what you lived and felt.
\end{grammarlens}

\subsection*{Your turn}
\begin{itemize}
\raggedright
  \item \emph{Say it:} \emph{akhlāq}
  \item \emph{Say it:} \emph{akhlāq}, once more
  \item \emph{Say it:} \emph{us ke akhlāq baṛe achchhe haiṅ.} — He has very good manners.
  \item \emph{Recall:} say the Urdu for courteous, then the Urdu for refinement, then say \emph{akhlāq} again
\end{itemize}

\begin{tcolorbox}[breakable,colback=teal!4,colframe=teal!35!black,title={Before you move on}]
What does \ur{اخلاق} mean? (\textquotedblleft{}Manners; moral conduct\textquotedblright{}.) Say it once more.
\end{tcolorbox}
\section[\texorpdfstring{\ur{احترام} \ur{کرنا} (ehtirām karnā)}{احترام کرنا (ehtirām karnā)}]{\ur{احترام} \ur{کرنا} (ehtirām karnā) — to show respect (to)}

\label{lesson:UR-C461-ehtiram-karna}



\begin{quote}
Before the new one: say the Urdu for refinement, then the Urdu for moral conduct.
\end{quote}

\subsection*{You'll want to know: \ur{احترام} \ur{کرنا}}
\textbf{\ur{احترام} \ur{کرنا}} — \emph{ehtirām karnā} — \textquotedblleft{}to show respect (to)\textquotedblright{}.

\begin{quote}
\ur{ہم} \ur{بزرگوں} \ur{کا} \ur{احترام} \ur{کرتے} \ur{ہیں۔}
\end{quote}

\begin{quote}
\emph{ham buzurgoṅ kā ehtirām karte haiṅ.}
\end{quote}

\begin{quote}
We show respect to our elders.
\end{quote}

Past, with \ur{نے} (he or she …): \ur{اس} \ur{نے} \ur{استاد} \ur{کا} \ur{احترام} \ur{کیا}, \emph{us ne ustād kā ehtirām kiyā}.

\begin{cousinweb}
\textbf{\ur{احترام} \ur{کرنا}} (\emph{ehtirām karnā}) is built from \textbf{\ur{احترام}}, usually traced, through Persian, to Arabic \emph{iḥtirām}, \textquotedblleft{}respect\textquotedblright{}, and \textbf{\ur{کرنا}}, \textquotedblleft{}to do\textquotedblright{}.
\end{cousinweb}

\begin{grammarlens}[title={What you've built}]
One more word for describing what you lived and felt.
\end{grammarlens}

\subsection*{Your turn}
\begin{itemize}
\raggedright
  \item \emph{Say it:} \emph{ehtirām karnā}
  \item \emph{Say it:} \emph{ehtirām karnā}, once more
  \item \emph{Say it:} \emph{ham buzurgoṅ kā ehtirām karte haiṅ.} — We show respect to our elders.
  \item \emph{Recall:} say the Urdu for refinement, then the Urdu for moral conduct, then say \emph{ehtirām karnā} again
\end{itemize}

\begin{tcolorbox}[breakable,colback=teal!4,colframe=teal!35!black,title={Before you move on}]
What does \ur{احترام} \ur{کرنا} mean? (\textquotedblleft{}To show respect (to)\textquotedblright{}.) Say it once more.
\end{tcolorbox}
